\section{Aplicación web}
\label{sec:web}

\subsection{Framework y justificación}

Se eligió React 18 con TypeScript en modo estricto sobre Vite, empaquetado en un
\emph{Dockerfile} \emph{multi-stage} (etapa \texttt{builder} con Node LTS, etapa de ejecución con
\texttt{nginx:alpine}). La aplicación sigue el patrón de \textbf{arquitectura por
características} (\emph{feature-based}): cada módulo de negocio (\texttt{console},
\texttt{admin}, \texttt{reports}, \texttt{auth}) agrupa sus propios componentes, \emph{hooks} de
acceso a la API y pruebas, en vez de una organización por tipo técnico (todos los componentes
juntos, todos los servicios juntos) que dificulta ubicar el código relacionado con una misma
funcionalidad de negocio.

\subsection{Capturas reales}
\label{sec:web-capturas}

La Figura~\ref{fig:web-capturas} muestra cuatro pantallas capturadas en vivo contra el
\emph{stack} completo en ejecución (no maquetas): inicio de sesión, la consola de operadores con
datos reales (incluido un ticket creado en vivo vía la API durante esta misma verificación, ver
\texttt{docs/evidencias/}), el panel de reportes (modelo de lectura CQRS de
\texttt{report-service}, con un total distinto al de la consola porque se reconstruye de forma
asíncrona a partir de eventos de Kafka, no es una copia sincrónica), y la administración de
cuentas.

\begin{figure}[H]
\centering
\begin{minipage}{0.48\textwidth}
  \centering
  \includegraphics[width=\linewidth]{../../release/screenshots/01_login.png}
  \\ \footnotesize (a) Inicio de sesión
\end{minipage}
\hfill
\begin{minipage}{0.48\textwidth}
  \centering
  \includegraphics[width=\linewidth]{../../release/screenshots/02_consola_operadores_admin.png}
  \\ \footnotesize (b) Consola de operadores (rol ADMIN)
\end{minipage}
\vspace{0.5em}
\begin{minipage}{0.48\textwidth}
  \centering
  \includegraphics[width=\linewidth]{../../release/screenshots/03_reportes.png}
  \\ \footnotesize (c) Reportes (CQRS)
\end{minipage}
\hfill
\begin{minipage}{0.48\textwidth}
  \centering
  \includegraphics[width=\linewidth]{../../release/screenshots/04_administracion.png}
  \\ \footnotesize (d) Administración de cuentas
\end{minipage}
\caption{Capturas reales de \texttt{apps/web} contra el sistema en ejecución (2026-09-03).}
\label{fig:web-capturas}
\end{figure}

\subsection{Rutas y control de acceso}

La aplicación cubre las cinco rutas mínimas exigidas por el Módulo B (\texttt{/}, \texttt{/login},
\texttt{/main}, \texttt{/settings}, \texttt{/about}), más dos adicionales con acceso restringido
por rol (\texttt{/admin}, \texttt{/reports}). El enrutamiento protegido se implementa con dos
componentes de orden superior: \texttt{ProtectedRoute} (exige sesión válida) y \texttt{RoleRoute}
(exige, además, un rol específico), ambos cubiertos por pruebas unitarias
(\texttt{ProtectedRoute.test.tsx}, \texttt{RoleRoute.test.tsx}).

El JSON Web Token~\cite{jones2015jwt} se guarda en \texttt{sessionStorage}, nunca en
\texttt{localStorage} sin cifrar: se pierde al cerrar la pestaña del navegador, acotando la
ventana de exposición ante un ataque de \emph{cross-site scripting} que lograra ejecutar código
en el origen de la aplicación.

\subsection{Internacionalización y estados explícitos}

La interfaz está completamente traducida a español e inglés con \texttt{react-i18next}, y
maneja de forma explícita los tres estados que toda pantalla que consume datos remotos debe
distinguir: carga (esqueletos animados en la tabla), error (mensaje con icono, nunca una pantalla
en blanco) y vacío (mensaje distinto a error cuando la consulta es válida pero no hay resultados).

La consola es \textbf{consciente del rol}: un usuario con rol \texttt{CLIENTE} ve un título y
subtítulo orientados a su propio caso de uso (``mis solicitudes'', no jerga interna de
operaciones), y en la columna de técnico asignado ve un distintivo genérico
(\texttt{AssignedChip}/\texttt{UnassignedChip}) en vez del identificador interno del técnico, que
no le aporta información accionable y expondría un identificador de otro usuario innecesariamente.

\subsection{Detalle completo del ticket}

Al hacer clic en cualquier fila de la tabla se abre un modal (\texttt{TicketDetailModal}) que
trae el recurso completo desde el backend (no solo los campos ya presentes en la fila),
incluyendo la línea de tiempo (creación, plazo de SLA, resolución si aplica) y el identificador
completo del ticket en un elemento seleccionable —pensado explícitamente para poder buscar ese
mismo identificador después, por ejemplo tras cerrarlo desde la aplicación móvil
(Sección~\ref{sec:movil}), y confirmar visualmente que ambos clientes leen el mismo estado del
mismo backend.

\subsection{Calidad}

\texttt{eslint} con las reglas por defecto de \texttt{@typescript-eslint} corre en modo estricto
(\texttt{--max-warnings 0}: cualquier advertencia rompe la compilación, no solo los errores). La
cobertura de pruebas (Vitest + Testing Library, \texttt{v8}) alcanza el 94.6\,\% de líneas sobre
el total de la aplicación —\texttt{TicketDetailModal.tsx}, que en una ronda anterior de esta
misma entrega era la brecha principal (34.8\,\%), ya está al 100\,\% de líneas hoy—, con la única
brecha real restante en \texttt{MouseParticles.tsx} (20.7\,\%): la estela de partículas
decorativa del fondo de \texttt{LoginPage}, canvas + \texttt{requestAnimationFrame} puramente
visual sin lógica de negocio, que no aporta cubrir exhaustivamente rama por rama.
